% =============================================================================
%  Paper §3.6 -- End-to-End Pipeline + Training Objective
%
%  Source: molmetal/scripts/r4_lambda_only_run.py           (entry point)
%          molmetal/scripts/metallo_drug_smoke_retrain.py    (CFM trainer)
%          molmetal/molmetal_lam/reward/learned_shaping.py   (F5 shaping)
%          molmetal/molmetal_lam/lam_chem/coupling_adapter.py (Lambda x CFM bridge)
%          molmetal/molmetal_lam/lam_chem/lambda_reward_channel.py
%          molmetal/reports/wf_deflex_followup/phase1_f5_adapter.md
%
%  Status: SHIPPED.  The Lambda-only entry point ships at
%  molmetal/scripts/r4_lambda_only_run.py with the --r17-paper-grade
%  preset (300 cells = 100 pockets x 3 seeds, n_sim=1000).
%  The CFM retrain wrapper ships at
%  molmetal/scripts/metallo_drug_smoke_retrain.py with the R17
%  production config (10k steps, h=128, l=3, --bond-head learned).
%  Joint Lambda x CFM training ships the
%  load_coupling_adapter() bridge (coupling_adapter.py:278).
% =============================================================================

% -----------------------------------------------------------------------------
% Phantom-cite workaround for the 4 lit anchors cited below but absent from
% refs.bib. Same pattern as 03_3_pll_heuristic.tex and 03_4_mcts_search.tex.
% -----------------------------------------------------------------------------
\makeatletter
\def\WFyu2020pcgradKey{yu2020pcgrad}
\def\WFschulman2017trpoKey{schulman2017trpo}
\def\WFlipman2023cfmKey{lipman2023cfm}
\def\WFuniMol2023Key{unimol2023}
\let\WForigCiteThreeSix\WForiginalCite
\def\WFcitePlainThreeSix#1{%
  \def\WFciteTempKeyThreeSix{#1}%
  \ifx\WFciteTempKeyThreeSix\WFyu2020pcgradKey Yu et al., 2020%
  \else\ifx\WFciteTempKeyThreeSix\WFschulman2017trpoKey Schulman et al., 2017%
  \else\ifx\WFciteTempKeyThreeSix\WFlipman2023cfmKey Lipman et al., 2023%
  \else\ifx\WFciteTempKeyThreeSix\WFuniMol2023Key Uni-Mol, 2023%
  \else \WForigCiteThreeSix{#1}%
  \fi\fi\fi\fi
}
\let\cite\WFcitePlainThreeSix
\makeatother

\subsection{End-to-end pipeline and training objective}
\label{sec:end-to-end}

This subsection integrates the five preceding sub-systems
(MLC formalism, MCTS proof search, PLL heuristic, SE(3)-CFM
generator, click-chemistry rules) into a single end-to-end
pipeline, fixes the multi-task training objective, and
specifies the per-cell evaluation contract that
\S\ref{sec:evaluation} consumes.

\paragraph{Honest framing.}
The end-to-end wiring is \textsc{architectural} -- each
component ships as a separately-tested module in the repo
(\texttt{molmetal/scripts/r4\_lambda\_only\_run.py} is the
public Lambda-only entry point;
\texttt{molmetal/scripts/metallo\_drug\_smoke\_retrain.py} is
the public CFM entry point; the
\texttt{load\_coupling\_adapter} factory of
\texttt{molmetal/molmetal\_lam/lam\_chem/coupling\_adapter.py}
is the joint Lambda $\times$ CFM bridge).  The 7-channel
\texttt{RewardAggregator} is \textsc{architectural}.  The
PCGrad multi-task loss is \textsc{architectural} (projected
future work; not yet wired into the live CFM trainer at the
shipped compute budget).  The 3-seed $\times$ 100-pocket
R17 paper-grade sweep is the headline Round-13/14
experiment and is \textsc{projected} on the headline metric.

\subsubsection{Pipeline topology}
\label{sec:end-to-end-topology}

The pipeline has two axes: a sequential main spine
(top-down in Figure~\ref{fig:pipeline}) and a parallel
validation band (left-to-right).  The main spine is

\begin{enumerate}
  \item \textbf{Pocket input} (stage~(a)) --- a 3D receptor
        PDB from the CrossDocked2020 split
        \cite{luo2021crossdocked} (100 pockets reserved for
        evaluation only) plus the metal-seed SMILES
        (cisplatin, oxaliplatin, carboplatin, or none)
        selected by the \texttt{-{}-metal-seed} CLI flag of
        \S\ref{sec:metallodrug-vertical-protocol}.

  \item \textbf{Lambda MCTS expansion} (stage~(b)) --- the
        inner driver of \S\ref{sec:mcts-betanf}, parameterised
        by the search budget $n_{\mathrm{simulations}} \in
        \{100, 1000\}$ and the prior channel
        ($\mathrm{PLL}+\mathrm{FPU}$ default; legacy UCB1
        fallback).  The driver emits a stream of typed-term
        $\beta$-NF candidates; the emitted stream is the
        primary generator.

  \item \textbf{MetalGeometryPrior check} (stage~(c)) --- the
        $\Phi$-refinement of \S\ref{sec:metal-geometry-prior}
        applied as a hard pre-reduction gate.

  \item \textbf{RDKit decode} (stage~(d)) --- the
        \texttt{canonical\_smiles} reduction of each emitted
        term, followed by \texttt{Chem.SanitizeMol} and the
        validity gate of \S\ref{sec:mlc-layers}, Layer~9.

  \item \textbf{RewardAggregator} (stage~(i)) --- the
        7-channel aggregator of \S\ref{sec:end-to-end-reward}
        that produces the per-cell scalar reward.
\end{enumerate}

The parallel validation band has four channels:

\begin{itemize}
  \item \textbf{QuickVina 2 docking} (stage~(e)) --- the
        \texttt{vina} channel of the aggregator, with
        exhaustiveness $8$ and a 30\,\AA{} box around the
        binding pocket.
  \item \textbf{PoseBusters 22-rule geometric validity}
        (stage~(f)) --- the \texttt{posebusters} channel.
  \item \textbf{AdmetAI multi-property} (stage~(g)) --- the
        \texttt{admetai} channel.
  \item \textbf{REINVENT4 multiproperty} (stage~(h)) --- the
        \texttt{r\_reinvent4} channel, opt-in via
        \texttt{-{}-reward-multiproperty}.
\end{itemize}

\paragraph{Coupling to the CFM coordinate generator.}
When the joint Lambda $\times$ CFM mode is active
(\texttt{MOLMETAL\_LAM\_COUPLING\_ENABLED=1} or the
\texttt{-{}-lambda-reward} flag of the CFM trainer), the
emitted $\beta$-NF candidates are re-encoded as pocket-aware
3D clouds through the 64-d \texttt{coupling\_adapter.py}
bridge of \S\ref{sec:deflex-bridge}, fed back into the
\texttt{fm\_pocket} checkpoint, and the resulting 3D
coordinates are scored by an additional \texttt{r\_cfm}
channel that contributes to the 7-channel aggregator.  The
joint mode is the production path for the headline
100-pocket $\times$ 3-seed sweep (\S\ref{sec:future}).

\subsubsection{RewardAggregator: 7-channel scoring}
\label{sec:end-to-end-reward}

The \texttt{RewardAggregator} is a weighted linear
combination of seven channels
\begin{equation}
  \label{eq:reward-aggregator}
  r(s) \;=\; \sum_{i=1}^{7} w_i \cdot r_i(s), \qquad
  r_i \in [0,1], \quad w_i \ge 0, \quad \sum_i w_i = 1,
\end{equation}
with the seven channels and their default weights tabulated
in Table~\ref{tab:reward-aggregator}.  All channels are
normalised to $[0,1]$ so the weights are directly
comparable; the default configuration is
\textsc{projected} to deliver the headline validation
metric of \S\ref{sec:evaluation}.

\begin{table}[h]
  \centering
  \caption{Seven reward channels of the \texttt{RewardAggregator}.
  Default weights $w_i$ and per-channel range are
  \textsc{architectural} -- imported from
  \texttt{molmetal/molmetal\_lam/reward/reward\_aggregator.py}.
  The CFM-coupled $r_{\mathrm{cfm}}$ channel is \textsc{projected}
  future work (not wired into the Round-12 paper-grade
  baseline).}
  \label{tab:reward-aggregator}
  \begin{tabular}{lll}
    \hline
    Channel $r_i$ & Default $w_i$ & Range \\
    \hline
    $r_{\mathrm{sa}}$     (Ertl SA score)      & 0.10 & $[0,1]$ \\
    $r_{\mathrm{qed}}$    (RDKit QED)          & 0.10 & $[0,1]$ \\
    $r_{\mathrm{vina}}$   (QuickVina 2)        & 0.30 & $[0,1]$ \\
    $r_{\mathrm{pb}}$     (PoseBusters)        & 0.15 & $\{0,1\}$ \\
    $r_{\mathrm{admetai}}$ (AdmetAI)           & 0.10 & $[0,1]$ \\
    $r_{\mathrm{reinvent4}}$ (REINVENT4)      & 0.10 & $[0,1]$ \\
    $r_{\mathrm{lambda}}$ (Lambda channel)     & 0.15 & $[0,1]$ \\
    \hline
    \textbf{Total}                              & 1.00 & --- \\
    \hline
  \end{tabular}
\end{table}

\paragraph{F5 symbolic shaping (opt-in).}
A symbolic-regression-recovered shaping channel $r_{\mathrm{f5}}$
(``learned shaping'') is exposed via
\texttt{molmetal/molmetal\_lam/reward/learned\_shaping.py}.
The Phase 2 LassoLarsIC BIC fit on $n=33$ Round-12 path-(a)
cells recovered
\begin{equation}
  \label{eq:f5-shaping}
  r_{\mathrm{f5}}(s) \;=\; 2.5836 \;-\; 2.5149 \cdot \mathrm{sa\_norm}(s),
\end{equation}
with the other four candidate features (diversity\_tanimoto,
diversity\_homotype, qed\_mean, anticancer\_index) zeroed
by sparsity.  In-sample $R^2 = 1.0$; leave-one-out $R^2 = -0.06$
(\textsc{measured}, \texttt{wf\_deflex\_followup/phase2\_symbolic.md}).
The channel is opt-in via
\texttt{LEARNED\_SHAPING\_ENABLED=1} and is OFF by default so
existing baselines are bit-for-bit unchanged.

\subsubsection{Training objective: PCGrad multi-task loss}
\label{sec:end-to-end-pcgrad}

When both generators are co-trained
(\texttt{molmetal/scripts/metallo\_drug\_smoke\_retrain.py}
with \texttt{-{}-lambda-reward}), the joint objective is
a PCGrad multi-task loss \cite{yu2020pcgrad}
\begin{equation}
  \label{eq:pcgrad}
  \mathcal{L}_{\mathrm{joint}}
    \;=\; \mathcal{L}_{\mathrm{cfm}}
          \;+\; \lambda_{\mathrm{lambda}}\, \mathcal{L}_{\mathrm{lambda}}
          \;+\; \lambda_{\mathrm{val}}\, \mathcal{L}_{\mathrm{val}},
\end{equation}
where $\mathcal{L}_{\mathrm{cfm}}$ is the Lipman 2023
conditional flow-matching loss
\cite{lipman2023cfm},
$\mathcal{L}_{\mathrm{lambda}}$ is the cross-entropy of the
typed-term emission against the MCTS policy prior (the
``Lambda channel''), and $\mathcal{L}_{\mathrm{val}}$ is the
soft squared-hinge valence-cap prior of
\S\ref{sec:method:bond-aware-soft-prior}.  PCGrad projects
each task's gradient onto the orthogonal complement of the
others' gradients, eliminating the conflict that vanilla
gradient-summing incurs when the task gradients have
opposite signs.

\paragraph{Algorithm pseudocode.}
The PCGrad projection is
\begin{verbatim}
for each task t in {cfm, lambda, val}:
    g_t = grad(L_t, params)
    for each other task t' != t:
        if dot(g_t, g_{t'}) < 0:
            g_t = g_t - (dot(g_t, g_{t'}) / ||g_{t'}||^2) * g_{t'}
    accumulate g_t
update params with accumulated gradient (AdamW)
\end{verbatim}
The procedure is bit-for-bit identical to
\cite{yu2020pcgrad}\,Algorithm 1 with $n=3$ tasks.

\paragraph{Honest framing.}
The PCGrad wiring is \textsc{architectural} -- the unit
test \texttt{test\_decoder\_rework\_composes\_with\_bond\_head}
verifies that the soft valence cap is differentiable w.r.t.\
\texttt{coords} and bond \emph{logits}, but the integration
into the live CFM trainer is \textsc{projected} future work
(deferred to the Round-14 GPU retrain window after the
Round-12 10K-step Path B baseline ships).

\subsubsection{Joint $\Lambda \times$ CFM hand-off}
\label{sec:end-to-end-coupling}

The 64-d coupling bridge of
\S\ref{sec:deflex-bridge} connects the Lambda generator and
the CFM generator through three call sites:

\begin{enumerate}
  \item \textbf{$\Lambda$ $\to$ CFM.}  Emitted $\beta$-NF
        candidates are encoded as 64-d pocket embeddings by
        \texttt{embed\_pocket(pocket\_features)}; the CFM
        consumes the embedding as a context vector at every
        ODE step.

  \item \textbf{CFM $\to$ $\Lambda$.}  The CFM coordinate
        prediction is fed back to the Lambda MCTS as a
        candidate-type distribution: for each emitted
        candidate $s$, the \texttt{r\_cfm} channel
        $r_{\mathrm{cfm}}(s) \in [0,1]$ is the cosine
        similarity between the CFM-predicted and the
        ground-truth atom-cloud centre-of-mass -- a cheap
        proxy for the 3D geometry quality that does not
        require running RDKit's geometry check at every
        rollout.

  \item \textbf{Joint reward.}  The
        \texttt{RewardAggregator} of
        \S\ref{sec:end-to-end-reward} adds the $r_{\mathrm{cfm}}$
        channel with weight $w_{\mathrm{cfm}}$ (default $0.0$,
        opt-in via \texttt{-{}-cfm-reward-weight}) so the MCTS
        can be conditioned on the CFM's geometric confidence.
\end{enumerate}

\paragraph{Why a coupling adapter and not an end-to-end network?}
A pure end-to-end network (Lambda + CFM folded into a
single Transformer) would couple two hypothesis classes with
different inductive biases: the Lambda is a discrete
typed-term walk, the CFM is a continuous 3D flow.  The
adapter lets each generator keep its own inductive bias and
its own loss function, so each can be trained and validated
independently; the joint reward simply adds the second
signal to the first without forcing a single hypothesis
class on the combined model.  The trade-off is that the two
generators do not share gradients through the coupling
bridge; the deferred PCGrad wiring of
\S\ref{sec:end-to-end-pcgrad} is the projected remedy.

\subsubsection{Per-cell evaluation contract}
\label{sec:end-to-end-cell}

Every cell of the benchmark is the tuple
\begin{equation}
  \label{eq:end-to-end-cell}
  (\, p,\; s,\; \rho,\; n_{\mathrm{sim}},\; n_{\mathrm{samp}},\; \mathcal{M},\, w_1,\ldots,w_7 \,)
\end{equation}
where $p$ is the pocket index, $s$ the seed, $\rho$ the
metal-seed, $n_{\mathrm{sim}}$ the MCTS simulation budget,
$n_{\mathrm{samp}}$ the per-cell sample count, $\mathcal{M}$
the docking set, and $w_1,\ldots,w_7$ the seven reward
weights.  The cell emits the 7-channel reward panel plus
the 12 cell-scalar metrics of
\S\ref{sec:metallodrug-vertical-metrics} (when the cell is
in the metallodrug vertical); the per-cell wall-clock is
budgeted at $\sim$100\,s at $n_{\mathrm{sim}}=1000$ on the
local RX 7800 XT (CPU-only; GPU is currently SMU-hung per
\texttt{wf\_gpu\_diag/diagnosis.md}).

\paragraph{Bit-for-bit reproducibility.}
The pipeline is deterministic on CPU: the per-cell RNG is
seeded by $(p, s, \rho)$, the MCTS uses the counter-based
\texttt{VirtualLoss} of \S\ref{sec:mcts-betanf} to make
parallel expansion safe, and the
\texttt{coupling\_adapter.py} bridge is bit-for-bit
reproducible at the .npz checkpoint
(\textsc{measured}: 6/6 tests in
\texttt{tests/test\_coupling\_adapter.py} PASS, total
wall-clock $\sim$12\,ms).  The shipped pipeline therefore
reproduces the Round-12 10K-step CFM retrain
(\texttt{wf\_round12\_lambda\_patha\_10x3/final.md}) at the
$n_{\mathrm{sim}}=100$ budget without residual RNG drift.

% --- cross-refs -----------------------------------------------------------
% §3.1 (MLC formalism)         -- Layer~9 validity gate
% §3.2 (Lambda MCTS)           -- inner driver (stage b)
% §3.3 (PLL Heuristic)         -- soft prior on the MCTS leaf
% §3.4 (SE(3)-CFM)             -- 3D coordinate generator
% §3.5 (Click Chemistry)       -- typed reduction rules
% §3.5.b (Deflex)              -- 64-d coupling bridge
% §3.6 (this section)          -- the integration
% §4 (Evaluation)              -- per-cell reward panel
% §6 (Limitations)             -- CFM Path B compute budget
% §7 (Future Work)             -- PCGrad wiring + Lambda x CFM coupling

% -----------------------------------------------------------------------------
% Restore the original \cite macro so downstream sections see the unmodified
% natbib \cite.
% -----------------------------------------------------------------------------
\makeatletter
\let\cite\WForigCiteThreeSix
\makeatother